\subsection{二次型的型.}

  给定 A 上向量空间 E 上的一个二次型 Q（§ 3 第 4 小节），我们称 E 为 Q 的定义空间，称 dim(E) 为 Q 的维数。给定 A 上向量空间 E、E' 上的两个二次型 Q、Q'，我们用 Q τ Q' 记它们的直和（§ 3 第 4 小节）。回忆，两个中性型的直和是中性的（§ 4 第 2 小节）。

  引入下列关系：

「Q 与 Q' 是 A 上有限维的非退化二次型，且存在中性二次型 N、N'，使得 Q τ N 等价于 Q' τ N'」。

  这个关系，我们将用 \(Q \sim Q'\) 来记，显然是自反的且对称的。它还是传递的：事实上，若 \(Q\)、\(Q'\)、\(Q''\) 是使 \(Q \sim Q'\) 与 \(Q' \sim Q''\) 成立的二次型，则存在中性二次型 \(M\)、\(M'\)、\(N\)、\(N'\)，使得 \(Q \tau M\) 等价于 \(Q' \tau M'\)，且 \(Q' \tau N\) 等价于 \(Q'' \tau N'\)；于是 \(Q \tau (M \tau N)\) 等价于 \((Q \tau M) \tau N\)，从而等价于 \((Q' \tau M') \tau N\)，也等价于 \((Q' \tau N) \tau M'\)，从而又等价于 \((Q'' \tau N') \tau M'\) 与 \(Q'' \tau (N' \tau M')\)；由于 \(M \tau N\) 与 \(N' \tau M'\) 是中性的，确实有 \(Q \sim Q''\)。因此，关系 \(Q \sim Q'\) 是 \(Q\) 与 \(Q'\) 之间的一个等价关系。显然，若 \(Q\) 与 \(Q'\) 是两个有限维且等价的非退化二次型，则有 \(Q \sim Q'\)。

  对 A 上每个非退化且有限维的二次型 Q，我们置

\[
\theta (Q) = \tau_ {x} (X \sim Q),\tag{1}
\]

并称 \(\theta(Q)\) 为 \(Q\) 的型。若 \(Q\) 与 \(Q'\) 是 \(A\) 上两个非退化且有限维的二次型，则关系 \(Q \sim Q'\) 与 \(\theta(Q) = \theta(Q')\) 等价。

命题 1. --- 设 Q 与 Q' 是 A 上两个非退化、有限维的二次型。要使 Q 与 Q' 等价，必须且只需它们具有相同的维数与相同的型。

  这个条件显然是必要的。设它被满足。于是存在中性型 N、N′，使得 Q ⊥ N 与 Q′ ⊥ N′ 等价。由于这两个型有相同的维数，N 与 N′ 亦然，从而它们等价（§ 4 第 2 小节命题 2 的推论 2）。于是由维特定理，Q 与 Q′ 等价（§ 4 第 3 小节定理 1 的推论 1）。

命题 2. --- 关系「存在 A 上一个非退化、有限维的二次型 Q，使得 X = θ(Q)」在 X 中是集合化的（《集合论》第二章 § 1 第 4 小节）。

  事实上，设 V 是 A 上一个无限维向量空间，S 是定义在 V 的有限维子空间上的非退化二次型所成的集合，W 是诸 \(\theta(Q)\)（\(Q \in \mathfrak{S}\)）所成的集合。显然，A 上每个非退化、有限维的二次型 \(Q'\) 等价于 \(\mathfrak{S}\) 的至少一个元素；由此 \(\theta(Q') \in \mathfrak{W}\)，这证明了我们的断言。

